\documentclass[10pt,pdf,utf8,english,compress,hyperref={unicode}]{beamer}
\usepackage{ifxetex}
\usepackage[english,russian]{babel}
\ifxetex
\else
\usepackage[utf8]{inputenc}
\usepackage[T2A]{fontenc}
\usepackage[backend=biber,style=authoryear]{biblatex}
\fi
\usepackage{booktabs}
\usepackage[scale=2]{ccicons}
\usepackage{url}
\usepackage{listings}
\usepackage{booktabs}
\usepackage{hyperref}
\usepackage{drawstack}

\graphicspath{{./pix/}}
%\usetheme[useTitleProgressBar]{m}
\usetheme{m}

% ------ Embed URL in the reference title ------

\ifxetex
% ---- START
\else
\renewbibmacro*{url+urldate}{}

\newbibmacro{string+doiurl}[1]{%
	\iffieldundef{doi}{%
		\iffieldundef{url}{%
			#1%
		}{%
			\href{\thefield{url}}{#1}%
		}%
	}{%
		\href{http://dx.doi.org/\thefield{doi}}{#1}%
	}%
}

\DeclareFieldFormat{title}{\usebibmacro{string+doiurl}{\mkbibemph{#1}}}
\DeclareFieldFormat[article,incollection]{title}{\usebibmacro{string+doiurl}{\mkbibemph{#1}}}

% ------ Nice and shiny bibliography list -------

\setbeamertemplate{bibliography item}{%
	\ifboolexpr{test {\ifentrytype{book}} or test {\ifentrytype{mvbook}}
	or test {\ifentrytype{collection}} or test {\ifentrytype{mvcollection}}
	or test {\ifentrytype{reference}} or test {\ifentrytype{mvreference}} }
	{\setbeamertemplate{bibliography item}[book]}
	{\ifentrytype{online}
		{\setbeamertemplate{bibliography item}[online]}
		{\setbeamertemplate{bibliography item}[article]}}%
	\usebeamertemplate{bibliography item}}

\defbibenvironment{bibliography}
	{\list{}
		{\settowidth{\labelwidth}{\usebeamertemplate{bibliography item}}%
		\setlength{\leftmargin}{\labelwidth}%
		\setlength{\labelsep}{\biblabelsep}%
		\addtolength{\leftmargin}{\labelsep}%
		\setlength{\itemsep}{\bibitemsep}%
		\setlength{\parsep}{\bibparsep}}}
	{\endlist}
	{\item}
% ----- END
\fi

% ------ Document starting here -------


\title{ESIL - Универсальный IL}
\subtitle{ESIL - Intermediate Language для radare2}
\author{Anton Kochkov (@akochkov)}
\date{\today}
\institute{ZeroNights 11-2015}

\ifxetex
\else
\addbibresource{refs.bib}
% Just a hack to workaround 'Package keyvar error: langjapanese undefined'
% Will be fixed in biblatex 3.2 and biber 2.3
% See more at:
% https://tex.stackexchange.com/questions/277314/package-keyval-error-langjapanese-undefined
\NewBibliographyString{langjapanese}
\NewBibliographyString{fromjapanese}
\fi

% ------------------------------------------------------------------------------------
\begin{document}
\maketitle

\begin{frame}[fragile]
  \frametitle{Anton Kochkov}
    \begin{itemize}
    \item Moscow, Russia
    \item Love Reverse Engineering, foreign languages and travel
	\item Member of R2 crew, radare2 evangelist
    \item Security Code Ltd.
    \end{itemize}
\end{frame}

% ------------------------------------------------------------------------------------
\section{Intermediate Languages}

\begin{frame}[fragile]
  \frametitle{What is Intermediate Language?}
  \begin{itemize}
  \item {\it Intermediate language is the language of an abstract machine designed to aid in the analysis
  of computer programs. \cite{il-wikipedia}}
  \item Heavily used for academic research and real world tools
  \item Vital for decompilation process
  \item Base for various kind of applications - SMT, AEG, AEP, etc
  \end{itemize}
\end{frame}

\begin{frame}[fragile]
  \frametitle{REIL \footfullcite{reil}}
  \begin{itemize}
  \item Invented by Zynamics company
  \item BinNavi, BinDiff were both based on top of REIL
  \item x86, ARM, PowerPC architectures supported
  \item Бесконечная память VM
  \item Бесконечное количество регистров VM
  \item Без Floating Point
  \item Оригинальные утилиты написаны на Java
  \end{itemize}
\end{frame}

\begin{frame}[fragile]
  \frametitle{REIL}
  \begin{itemize}
  \item 17 инструкций
  \item Алиасы для реальных регистров (eax, ebx, r0, \ldots) \footfullcite{reil-zynamics}
  \end{itemize}
  % TODO: Make a table of them
\end{frame}

\begin{frame}[fragile]
  \frametitle{BAP}
  \begin{itemize}
  \item BAP - Binary Analysis Platform \footfullcite{bap-handbook}
  \item Настоящее имя IL - BIL
  \item Развитый фреймворк
  \item Интеграция с другими утилитами - TEMU, libVEX, IDA Pro, Qira, \ldots
  \item Ориентирован на x86, ARM
  \item Без Floating Point
  \end{itemize}
\end{frame}

%
%\begin{frame}[fragile]
%  \frametitle{BAP}
%  \begin{itemize}
%  \item DEMO
%  \end{itemize}
%\end{frame}

\begin{frame}[fragile]
  \frametitle{BitBlaze (VineIL/VEX)}
  \begin{itemize}
  \item BitBlaze \footfullcite{bitblaze} - платформа, аналогичная BAP
  \item Имеет несколько промежуточных языков
  \item VEX IL (libVEX из valgrind) - ``нижний'' уровень
  \item Vine IL - ``верхний'' уровень
  \item Написан на OCaml + C++
  \end{itemize}
\end{frame}

\begin{frame}[fragile]
  \frametitle{VEX IL}
  \begin{itemize}
  \item Явное указание всех side-эффектов для команд
  \item Ближе всего к ESIL
  \item Оттестирован и используется в Valgrind
  \item Хорошо подходит для эмуляции кода
  \item Избыточен
  \end{itemize}
\end{frame}

\begin{frame}[fragile]
  \frametitle{Vine IL \footfullcite{vine-manual} \footfullcite{vine-vuln-analysis}}
  \begin{itemize}
  \item Бесконечная память
  \item Бесконечное количество регистров
  \item Поддержка типов
  \item Поддержка ``variable scope''
  \end{itemize}
\end{frame}

\begin{frame}[fragile]
  \frametitle{RREIL, OpenREIL, MAIL}
  \begin{itemize}
  \item RREIL \footfullcite{rreil} - гибкий язык, замена REIL
  \item RREIL - поддержка типов
  \item RREIL - интересная концепция ``доменов''
  \item MAIL - IL, созданный для анализа Malware
  \item MAIL - позволяет программе перезаписывать себя саму
  \item RREIL и MAIL - опять отсутствие Floating Point
  \end{itemize}
\end{frame}

\begin{frame}[fragile]
  \frametitle{RREIL, OpenREIL, MAIL}
  \begin{itemize}
  \item OpenREIL \footfullcite{openreil} - проект, созданный для использования REIL в современных реалиях
  \item OpenREIL - полноценный фреймворк, как и BAP
  \item OpenREIL отличается от оригинального REIL
  \item Использует libVEX и имеет поддержку SMT-solving
  \end{itemize}
\end{frame}


% -------------------------------------------------------------------------------------
\section{ESIL - сходства и отличия}

\begin{frame}[fragile]
  \frametitle{Краткое описание}
     \begin{itemize}
        \item Evaluable Strings Intermediate Language \footfullcite{r2esil}
		\item Использует обратную польскую нотацию (для скорости)
		\item Не предназначен для чтения человеком
		\item По ``уровню'' приближен к VEX
		\item Небольшое количество инструкций
		\item Полный учёт side-эффектов
      \end{itemize}
\end{frame}

\begin{frame}[fragile]
  \frametitle{Краткое описание}
     \begin{itemize}
        \item Сроектирован для большого количества архитектур
		\item Бесконечная память
		\item Бесконечные регистры
		\item Алиасы (использование ``нативных'' имен)
		\item Есть возможность вызывать куски нативного кода (+syscall)
		\item Возможность добавления ``custom ops''
		\item Нет Floating Point (будет в следующей версии)
      \end{itemize}
\end{frame}

\begin{frame}[fragile]
  \frametitle{Операнды ESIL}
    \begin{table}[H]
	\caption{ESIL Operands \footfullcite{r2esil-instruction-set}}
	\begin{center}
		\resizebox{\linewidth}{!}{%
		\begin{tabular}{|*{4}{c|}}
		\hline
		\textbf{ESIL Opcode} & Operands & Name & Desription\\
		\hline
		\$ & src & Syscall & syscall \\
		\hline
		\$\$ & src & Instruction address & Get address of current instruction \\
		\hline
		== & src,dst & Compare & v = dst - src ; update\_eflags(v) \\
		\hline
		\textless &	src,dst & Smaller &	stack = (dst \textless src) \\
		\hline
		\textless= & src,dst & Smaller or Equal	& stack = (dst \textless= src) \\
		\hline
		\textgreater &	src,dst & Bigger & stack = (dst \textgreater src) \\
		\hline
		|= & src,reg & OR eq & reg = reg | src \\
		\hline
		\end{tabular}}
	\end{center}
	\end{table}
\end{frame}

% -------------------------------------------------------------------------------------
\section{Практическое применение}

\begin{frame}[fragile]
	\frametitle{Radare}
	\begin{figure}
		\includegraphics[scale=0.2]{r2-ubahn.jpg}\footfullcite{r2advert}
	\end{figure}
\end{frame}

\begin{frame}[fragile]
  \frametitle{Radare2 утилиты}
     \begin{itemize}
        \item rax2
        \item rabin2
        \item rasm2
        \item radiff2
        \item rafind2
        \item rahash2
        \item radare2
        \item r2pm
        \item rarun2/ragg2/ragg2-cc
      \end{itemize}
\end{frame}

\begin{frame}[fragile]
  \frametitle{1 command \textless—\textgreater 1 Reverse-Engineering'notion}
  \begin{enumerate}
  \item Каждый символ команды что-то значит \alert{(w = write, p = print)}
  \item Обычно команды - это аббревиатуры действий \alert{pdf = p \textless-\textgreater print d \textless-\textgreater disassemble f \textless-\textgreater function }
  \item Доступна короткая справка для каждой команды \textbf{cmd?}, например \alert{pdf?},\alert{?}, \alert{???}, \alert{???}, \alert{?\$?}, \alert{?@?}
  \end{enumerate}
\end{frame}

\begin{frame}[fragile]
  \frametitle{Radare2 — Основные команды CLI-режима}
  \begin{enumerate}
   \item \alert{\bf{r2 -A}} или \alert{\bf{r2}} + \alert{\bf{aaa}} : Анализ
   \item \alert{\bf{s}} : Переход по указанному адресу
   \item \alert{\bf{pdf}} : Дизассемблирование функции
   \item \alert{\bf{af?}} : Анализ функции
   \item \alert{\bf{ax?}} : Анализ XREF
   \item \alert{\bf{/?}} : Поиск
   \item \alert{\bf{ps?}} : Напечатать строку (print string)
   \item \alert{\bf{C?}} : Комментарии
   \item \alert{\bf{w?}} : Запись (hex, опкодов, etc)
 \end{enumerate}
\end{frame}

\section{Radare2 — Visual mode}
\begin{frame}[fragile]
  \frametitle{Radare2 — Основные команды визуального режима}
  \begin{enumerate}
  \item \alert{\bf{V?}} или просто \alert{\bf{?}} : Помощь по командам
  \item \alert{\bf{p/P}} : переключение между разными визуальными представлениями
  \item Навигация с помощью стрелок/hjkl
  \item \alert{\bf{o}} : переместиться по адресу
  \item \alert{\bf{e}} : визуальный режим настроек
  \item \alert{\bf{v}} : список функций
  \item \alert{\bf{\_}} : HUD
  \item \alert{\bf{V}} : ASCII Graph
  \item \alert{\bf{0-9}} : Прыжок на функцию
  \item \alert{\bf{u}} : Undo
 \end{enumerate}
\end{frame}

\begin{frame}[fragile]
  \frametitle{Эмуляция участков кода}
     \begin{itemize}
        \item \alert{\bf{ae*}} набор инструкций
		\item \alert{\bf{aei}} - инициализация ESIL VM
		\item \alert{\bf{aeim}} - инициализация стека/памяти VM
		\item \alert{\bf{aeip}} - установка IP (Instruction Pointer)
		\item \alert{\bf{aes}} - step в режиме эмуляции ESIL
		\item \alert{\bf{aec[u]}} - continue [until]
		\item \alert{\bf{aef}} - эмуляция функции
      \end{itemize}
\end{frame}

\begin{frame}[fragile]
  \frametitle{Эмуляция участков кода}
     \begin{itemize}
        \item DEMO
      \end{itemize}
\end{frame}

% show demo here - demos/demo5_it8502e
\begin{frame}[fragile]
\ifxetex
  \frametitle{Embedded controller - 8051 - ESIL VM}
\else
  \frametitle{Embedded controller - 8051 - ESIL VM \footfullcite{r2esiltv}}
\fi
  \begin{itemize}
	  \item \alert{\bf{r2 -a 8051 ite\_it8502.rom}}
	  \item \alert{\bf{. ite\_it8502.r2}}
	  \item \alert{\bf{e io.cache=true}} для использования кеширования IO
	  \item запустим \alert{\bf{aei}}
	  \item запустим \alert{\bf{aeim}}
	  \item запустим \alert{\bf{aeip}} для старта с момента указания команды
	  \item \alert{\bf{aecu [addr]}} для эмуляции, пока не достигнем IP = [addr]
  \end{itemize}
\end{frame}
% show demo here
% Reference to the example from here? http://radare.tv/a/44

% Check before the show! Probably won't work
\begin{frame}[fragile]
  \frametitle{Совместная отладка}
     \begin{itemize}
		\item Использование ``подсказок'' ESIL при визуальной отладке
        \item DEMO
      \end{itemize}
\end{frame}

\begin{frame}[fragile]
  \frametitle{Эмуляция VM}
     \begin{itemize}
        \item Позволяет выполнить распаковку или выполнение в VM
		\item Хороший пример - использование ESIL для распаковки Baleful \footfullcite{esil-baleful}
      \end{itemize}
\end{frame}

\begin{frame}[fragile]
  \frametitle{Автоматическое отображение результатов эмуляции в дизассемблере}
     \begin{itemize}
        \item Отображает в комментариях значения регистров и пямяти
		\item Использует тот же механизм эмуляции кода ESIL VM
		\item Показывает likely/unlikely для условных переходов
		\item \alert{\bf{e asm.emu=true}}
      \end{itemize}
\end{frame}

\begin{frame}[fragile]
  \frametitle{Автоматическое отображение результатов эмуляции в дизассемблере}
     \begin{itemize}
        \item DEMO
      \end{itemize}
\end{frame}

\begin{frame}[fragile]
  \frametitle{Конвертация в другие языки - OpenREIL}
     \begin{itemize}
        \item OpenREIL - развитый фреймфорк
		\item Есть возможность использования SMT
		\item Добавлена возможность конфертации ESIL в OpenREIL
		\item Команда \alert{\bf{aetr}}
      \end{itemize}
\end{frame}

\begin{frame}[fragile]
  \frametitle{Конвертация в другие языки - OpenREIL}
     \begin{itemize}
        \item DEMO
      \end{itemize}
\end{frame}

\begin{frame}[fragile]
  \frametitle{Embedded controller - 8051 - ESIL2REIL}
  \begin{itemize}
	  \item \alert{\bf{r2 -a 8051 ite\_it8502.rom}}
	  \item \alert{\bf{. ite\_it8502.r2}}
	  \item run \alert{\bf{pae 36}} для показа ESIL представления функции 'set\_SMBus\_frequency'
\ifxetex
	  \item \alert{\bf{aetr \`{}pae 36\`}} для конвертации строки ESIL в REIL
\else
	  \item run aetr \`{}pae 36\`{} для конвертации строки ESIL в REIL\footfullcite{openreil}
\fi
	  \item Сохранить вывод в файл и передать управление в OpenREIL
	  \item Можно проделать всё вышеперечисленное с помощью r2pipe скрипта
  \end{itemize}
\end{frame}

% --------------------------------------------------------------------------------------
\section{Radeco IL и radeco decompiler}

\begin{frame}[fragile]
  \frametitle{ESIL -> Radeco \footfullcite{radeco}}
     \begin{itemize}
        \item Использует ESIL в качестве входных данных
		\item Использует другие метаданные из radare2
		\item Соединяется с radare2 через r2pipe
		\item Написан на Rust
		\item Большая часть кода написана двумя студентами GSoC 2015
		\item Авторы - Sushant Dinesh и David Kreuter \footfullcite{radeco-gsoc}
		\item GSoC 2015 прошел под эгидой проекта Openwall
      \end{itemize}
\end{frame}

\begin{frame}[fragile]
  \frametitle{Причины появления декомпилятора}
     \begin{itemize}
        \item Существующие FOSS декомпиляторы не учитывают последние исследования
		\item Академические (но интересные) идеи не имеют полноценной реализации
		\item Radare2 нуждается в декомпиляторе
		\item Хорошее и интересное задание для Google Summer of Code
      \end{itemize}
\end{frame}

\begin{frame}[fragile]
  \frametitle{Описание Radeco IL}
     \begin{itemize}
        \item Графовое представление
		\item Взяты идеи из RREIL и MAIL
		\item Использование SSA на этапе лифтинга ESIL -> Radeco IL
		\item Встроенная поддержка DCE (Dead Code Elimination)
		\item Базовая возможность вывода типов \footfullcite{tie}
      \end{itemize}
\end{frame}

\begin{frame}[fragile]
  \frametitle{Radeco demo}
     \begin{itemize}
        \item DEMO
      \end{itemize}
\end{frame}

% --------------------------------------------------------------------------------------
\section{Пути будущего развития}

\begin{frame}[fragile]
  \frametitle{Поддерживаемые архитектуры}
     \begin{itemize}
        \item Сейчас лучше всего поддерживаются x86, ARM, GameBoy, 8051, etc
		\item Глобальная цель - поддержка ESIL для всех архитектур в radare2
		\item Поддержка профилей для выбранных модификаций/вариаций процессоров
      \end{itemize}
\end{frame}

\begin{frame}[fragile]
  \frametitle{Поддерживаемые наборы инструкций}
     \begin{itemize}
        \item Floating point (LLVM/McSema) \footfullcite{floating-point-SO}
		\item Векторные инструкции (SSE, AVX, Neon, etc)
		\item VLIW инструкции (для эмуляции кода DSP)
      \end{itemize}
\end{frame}

\begin{frame}[fragile]
  \frametitle{Визуальная отладка и трассировка}
     \begin{itemize}
        \item Улучшение UI
		\item Возможность визуального сравнения эмуляции и нативного выполнения
		\item Устранение ``мертвого'' кода из ASCII графов на лету
		\item Интеграция в WebUI и Bokken \footfullcite{bokken}
      \end{itemize}
\end{frame}

\begin{frame}[fragile]
  \frametitle{Развитие декомпилятора Radeco}
     \begin{itemize}
        \item Генерация C кода
		\item Поддержка нативных типов
		\item Синхронизация с отладкой
		\item Автовывод типов/распознавание объектов и классов\footfullcite{tie} \footfullcite{il-optimization}
      \end{itemize}
\end{frame}

% ----------------------------------------------------------------------------------------
\ifxetex
\else
\section{References}
\begin{frame}[allowframebreaks]
	\frametitle{A lot of them}
	\printbibliography
\end{frame}
\fi

\end{document}
